\documentclass[12pt]{jlreq}
\usepackage{amsmath, amssymb}

\newcommand{\C}{\mathbb{C}}
\newcommand{\R}{\mathbb{R}}
\newcommand{\Z}{\mathbb{Z}}
\newcommand{\N}{\mathbb{N}}
\newcommand{\F}{\mathbb{F}}
\newcommand{\Prob}{\mathbb{P}}
\newcommand{\Exp}{\mathbb{E}}
\newcommand{\dd}{\,d}
\newcommand{\Ker}{\operatorname{Ker}}
\newcommand{\Impart}{\operatorname{Im}}
\newcommand{\Hom}{\operatorname{Hom}}

\begin{document}

\(X\) を \(\R^3\) から原点 \(O\) を除いた領域とする．\(f(x, y, z)\) は \(X\) 上の \(C^\infty\) 級関数で，\(r = \sqrt{x^2+y^2+z^2}\) を用いて，\(f(x, y, z) = h(r)\) と表されるとする．\(X\) 上の \(1\) 次微分形式 \(\omega\) を
\[
  \omega = f(x, y, z) (x \, dx + y \, dy + z \, dz)
\]
で定める．

(1) \(\omega\) は \(X\) 上の閉形式であること，つまり，\(X\) 上 \(d\omega = 0\) が成立することを示せ．

(2) \(\omega\) は \(X\) 上の完全形式であること，つまり，\(X\) 上のある \(C^\infty\) 級関数 \(g\) によって，\(\omega = dg\) と表されることを示せ．

(3) \(\omega\) が \(\Delta \phi = 0\) を満たす \(X\) 上の \(C^\infty\) 級関数 \(\phi\) によって，\(\omega = d\phi\) となるとき，関数 \(f\) を \(x, y, z\) を用いて表せ．ここで，
\[
  \Delta \phi = \frac{\partial^2 \phi}{\partial x^2} + \frac{\partial^2 \phi}{\partial y^2} + \frac{\partial^2 \phi}{\partial z^2}
\]
である．

\end{document}